%document
\documentclass[10pt]{beamer}
%theme
\usetheme{metropolis}
% packages
\usepackage{color}
\usepackage{listings}
\usepackage[ngerman]{babel}
\usepackage[utf8]{inputenc}
\usepackage{multicol}

\usepackage{hyperref}
\hypersetup{colorlinks=true, linkcolor=blue, urlcolor=red}

% color definitions
\definecolor{mygreen}{rgb}{0,0.6,0}
\definecolor{mygray}{rgb}{0.5,0.5,0.5}
\definecolor{mymauve}{rgb}{0.58,0,0.82}

\lstset{
    backgroundcolor=\color{white},
    % choose the background color;
    % you must add \usepackage{color} or \usepackage{xcolor}
    basicstyle=\footnotesize\ttfamily,
    % the size of the fonts that are used for the code
    breakatwhitespace=false,
    % sets if automatic breaks should only happen at whitespace
    breaklines=true,                 % sets automatic line breaking
    captionpos=b,                    % sets the caption-position to bottom
    commentstyle=\color{mygreen},    % comment style
    % deletekeywords={...},
    % if you want to delete keywords from the given language
    extendedchars=true,
    % lets you use non-ASCII characters;
    % for 8-bits encodings only, does not work with UTF-8
    frame=single,                    % adds a frame around the code
    keepspaces=true,
    % keeps spaces in text,
    % useful for keeping indentation of code
    % (possibly needs columns=flexible)
    keywordstyle=\color{blue},       % keyword style
    % morekeywords={*,...},
    % if you want to add more keywords to the set
    numbers=left,
    % where to put the line-numbers; possible values are (none, left, right)
    numbersep=5pt,
    % how far the line-numbers are from the code
    numberstyle=\tiny\color{mygray},
    % the style that is used for the line-numbers
    rulecolor=\color{black},
    % if not set, the frame-color may be changed on line-breaks
    % within not-black text (e.g. comments (green here))
    stepnumber=1,
    % the step between two line-numbers.
    % If it's 1, each line will be numbered
    stringstyle=\color{mymauve},     % string literal style
    tabsize=4,                       % sets default tabsize to 4 spaces
    % show the filename of files included with \lstinputlisting;
    % also try caption instead of title
    language = C,
	showspaces = false,
	showtabs = false,
	showstringspaces = false,
	escapechar = ,
}

\def\ContinueLineNumber{\lstset{firstnumber=last}}
\def\StartLineAt#1{\lstset{firstnumber=#1}}
\let\numberLineAt\StartLineAt



\newcommand{\codeline}[1]{
	\alert{\texttt{#1}}
}

% This Document contains the information about this course.

% Authors of the slides
\author{
    Lecturers:
	%name1
	Mirko Jantschke,
	%name2
	Pascal Scholz
	\\
	%creaters of the slides
	Created by: Richard Mörbitz, Manuel Thieme
}


% Fancy Logo
\titlegraphic{\hfill\includegraphics[height=1.25cm]{../templates/fsr_logo_cropped}}


\title{C-Lessons}
\subtitle{Variables}
\date{\today}



% the actual document
\begin{document}

\maketitle

%-------------------------------------------------------------------------------------------------
\begin{frame}{Contents}
	\tableofcontents
\end{frame}

%-------------------------------------------------------------------------------------------------
\section{Lists and binary trees}
%-------------------------------------------------------------------------------------------------

\begin{frame}[fragile]{Yay, pointers}
	For better handling of datastructures, we define \textit{structs} with additional pointers to themselves:\\\ \\
	\begin{columns}[c]
		\column{.5\textwidth}
		\begin{lstlisting}[numbers=none]
struct listElement {
	int value;
	struct listElement *next;
}
\end{lstlisting}
		\column{.5\textwidth}
		\begin{tikzpicture}[scale=1.5,font=\scriptsize]
			\node[font=\small] at (0,1) [above] {struct listElement};
			
			\draw (0,0) -- (0,1);
			\draw (0,0) -- (2,0);
			\draw (0,1) -- (2,1);
			\draw (2,0) -- (2,1);
			
			\draw (.1,.1) -- (.1,.45);
			\draw (.1,.1) -- (1.9,.1);
			\draw (.1,.45) -- (1.9,.45);
			\draw (1.9,.1) -- (1.9,.45);
			\node at (0.1,.275) [right] {struct listElement*};
			
			\node[circle, thick, draw=black, minimum size=10pt, inner sep = 0pt] at (1.72,.275) {};
			\draw [thick,->,shorten >=.2ex] (1.85,.257) to[out=0,in=-90] (2.15,1) to[out=90,in=0] (1.5,1.25) to[out=180,in=90](1,1);
			
			\draw (.1,.55) -- (.1,.9);
			\draw (.1,.55) -- (1.9,.55);
			\draw (.1,.9) -- (1.9,.9);
			\draw (1.9,.55) -- (1.9,.9);
			\node at (0.1,.725) [right] {int};
		\end{tikzpicture}
	\end{columns}
\end{frame}

%-------------------------------------------------------------------------------------------------

\begin{frame}{Single-linked lists}

	\begin{tikzpicture}[scale=.73,font=\scriptsize]
			
		\draw (-2.5,.1) -- (-1,.1);
		\draw (-2.5,.1) -- (-2.5,.45);
		\draw (-2.5,.45) -- (-1,.45);
		\draw (-1,.45) -- (-1,.1);
		\node[circle, thick, draw=black, minimum size=5pt, inner sep = 0pt] at (-1.18,.275) {};
		\node at (-2.6,.275) [right] {head};
		\draw (-1.05,.275) edge[thick,out=0,in=135,->,shorten >=.2ex] (0,1);
	
		\draw (0,0) -- (0,1);
		\draw (0,0) -- (2,0);
		\draw (0,1) -- (2,1);
		\draw (2,0) -- (2,1);
		
		\draw (.1,.1) -- (.1,.45);
		\draw (.1,.1) -- (1.9,.1);
		\draw (.1,.45) -- (1.9,.45);
		\draw (1.9,.1) -- (1.9,.45);
		\node at (0,.275) [right] {next};
		
		\node[circle, thick, draw=black, minimum size=5pt, inner sep = 0pt] at (1.72,.275) {};
		\begin{uncoverenv}<-3>
			\draw (1.85,.275) edge[thick,out=0,in=135,->,shorten >=.2ex] (3,1);
		\end{uncoverenv}
		\begin{uncoverenv}<4>
			\draw (1.85,.275) edge[red,thick,out=0,in=135,->,shorten >=.2ex] (3,1);
		\end{uncoverenv}
		
		\draw (.1,.55) -- (.1,.9);
		\draw (.1,.55) -- (1.9,.55);
		\draw (.1,.9) -- (1.9,.9);
		\draw (1.9,.55) -- (1.9,.9);
		\node at (0,.725) [right] {value = ...};
		
		
		\draw (3,0) -- (3,1);
		\draw (3,0) -- (5,0);
		\draw (3,1) -- (5,1);
		\draw (5,0) -- (5,1);
		
		\draw (3.1,.1) -- (3.1,.45);
		\draw (3.1,.1) -- (4.9,.1);
		\draw (3.1,.45) -- (4.9,.45);
		\draw (4.9,.1) -- (4.9,.45);
		\node at (3,.275) [right] {next};
		
		\node[circle, thick, draw=black, minimum size=5pt, inner sep = 0pt] at (4.72,.275) {};
		\draw (4.85,.275) edge[thick,out=0,in=135,->,shorten >=.2ex] (6,1);
		
		
		\draw (3.1,.55) -- (3.1,.9);
		\draw (3.1,.55) -- (4.9,.55);
		\draw (3.1,.9) -- (4.9,.9);
		\draw (4.9,.55) -- (4.9,.9);
		\node at (3,.725) [right] {value = ...};
		
		
		\draw (6,0) -- (6,1);
		\draw (6,0) -- (8,0);
		\draw (6,1) -- (8,1);
		\draw (8,0) -- (8,1);
		
		\draw (6.1,.1) -- (6.1,.45);
		\draw (6.1,.1) -- (7.9,.1);
		\draw (6.1,.45) -- (7.9,.45);
		\draw (7.9,.1) -- (7.9,.45);
		\node at (6,.275) [right] {next};
		
		\node[circle, thick, draw=black, minimum size=5pt, inner sep = 0pt] at (7.72,.275) {};
		\draw (7.85,.275) edge[thick,out=0,in=135,->,shorten >=.2ex] (9,1);
		
		\draw (6.1,.55) -- (6.1,.9);
		\draw (6.1,.55) -- (7.9,.55);
		\draw (6.1,.9) -- (7.9,.9);
		\draw (7.9,.55) -- (7.9,.9);
		\node at (6,.725) [right] {value = ...};
		
		
		\draw (9,0) -- (9,1);
		\draw (9,0) -- (11,0);
		\draw (9,1) -- (11,1);
		\draw (11,0) -- (11,1);
		
		\draw (9.1,.1) -- (9.1,.45);
		\draw (9.1,.1) -- (10.9,.1);
		\draw (9.1,.45) -- (10.9,.45);
		\draw (10.9,.1) -- (10.9,.45);
		\node at (9,.275) [right] {next};
		
		
		\node[circle, thick, draw=black, minimum size=5pt, inner sep = 0pt] at (10.72,.275) {};
		\draw (10.85,.275) edge[thick,out=0,in=155,->,shorten >=.2ex] (11.5,.45);
		
		\draw (9.1,.55) -- (9.1,.9);
		\draw (9.1,.55) -- (10.9,.55);
		\draw (9.1,.9) -- (10.9,.9);
		\draw (10.9,.55) -- (10.9,.9);
		\node at (9,.725) [right] {value = ...};

		\draw (11.5,.1) -- (11.5,.45);
		\draw (11.5,.1) -- (12.5,.1);
		\draw (11.5,.45) -- (12.5,.45);
		\draw (12.5,.1) -- (12.5,.45);
		\node at (11.4,.265) [right] {NULL};
		
		
		\begin{uncoverenv}<2->
			\draw (2,-.5) -- (2,-1.5);
			\draw (2,-.5) -- (4,-.5);
			\draw (4,-.5) -- (4,-1.5);
			\draw (4,-1.5) -- (2,-1.5);
			
			\draw (2.1,-.6) -- (2.1, -.95);
			\draw (2.1,-.6) -- (3.9,-.6);
			\draw (3.9,-.6) -- (3.9,-.95);
			\draw (3.9,-.95) -- (2.1,-.95);
			\node at (2,-.775) [right] {value = ...};
			
			
			\draw (2.1, -1.05) -- (2.1, -1.4);
			\draw (2.1,-1.05) -- (3.9,-1.05);
			\draw (3.9,-1.05) -- (3.9,-1.4);
			\draw (3.9,-1.4) -- (2.1,-1.4);
			\node at (2,-1.225) [right] {next};
			\node[circle, thick, draw=black, minimum size=5pt, inner sep = 0pt] at (3.72,-1.225) {};
			
			
		\end{uncoverenv}
		\begin{uncoverenv}<3->
			\draw [thick,blue,->,shorten >=.2ex] (3.85,-1.225) to[out=0,in=-35] (4,-.4) to[out=180-35,in=-45] (2.5,-.1) to[out=180-45,in=135] (3,1);
		\end{uncoverenv}
		\begin{uncoverenv}<5->
			\draw [thick,red,->,shorten >=.2ex] (1.72,.145) to[out=-90,in=135] (2,-.5);
		\end{uncoverenv}
		
		
	\end{tikzpicture}
	\ \\\ \\
	\begin{itemize}[<+->]
		\item Now we can build lists of this structures.
		\item If we want to add a value at a specific position...
		\item ... we just have to change the pointers.
	\end{itemize}
	
	\begin{itemize}
		\item<6-> Additionally we can easily iterate through this list.
		\item<6-> But only in one direction.
	\end{itemize}
	
\end{frame}

%-------------------------------------------------------------------------------------------------

\begin{frame}{Suddenly: Double-linked lists}

	\begin{tikzpicture}[scale=.73,font=\scriptsize]
			
		\draw (-2.5,.1) -- (-1,.1);
		\draw (-2.5,.1) -- (-2.5,.45);
		\draw (-2.5,.45) -- (-1,.45);
		\draw (-1,.45) -- (-1,.1);
		\node[circle, thick, draw=black, minimum size=5pt, inner sep = 0pt] at (-1.18,.275) {};
		\node at (-2.6,.275) [right] {head};
		\draw (-1.05,.275) edge[thick,out=0,in=135,->,shorten >=.2ex] (0,1);
		
		\draw (-2,-.1) -- (-1,-.1);
		\draw (-2,-.1) -- (-2,-.45);
		\draw (-2,-.45) -- (-1,-.45);
		\draw (-1,-.45) -- (-1,-.1);
		\node at (-2.1,-.275) [right] {NULL};
	
		\draw (0,-.45) -- (0,1);
		\draw (0,-.45) -- (2,-.45);
		\draw (0,1) -- (2,1);
		\draw (2,-.45) -- (2,1);
		
		\draw (.1,-.35) -- (.1,0);
		\draw (.1,-.35) -- (1.9,-.35);
		\draw (.1,0) -- (1.9,0);
		\draw (1.9,-.35) -- (1.9,0);
		\node at (.3,-.225) [right] {prev};
		
		\node[circle, thick, draw=black, minimum size=5pt, inner sep = 0pt] at (.28,-.175) {};
		\draw (.15,-.175) edge[thick,out=180,in=-35,->,shorten >=.2ex] (-1,-.45);
		
		\draw (.1,.1) -- (.1,.45);
		\draw (.1,.1) -- (1.9,.1);
		\draw (.1,.45) -- (1.9,.45);
		\draw (1.9,.1) -- (1.9,.45);
		\node at (0,.275) [right] {next};
		
		\node[circle, thick, draw=black, minimum size=5pt, inner sep = 0pt] at (1.72,.275) {};
		\begin{uncoverenv}<-3>
			\draw (1.85,.275) edge[thick,out=0,in=135,->,shorten >=.2ex] (3,1);
		\end{uncoverenv}
		\begin{uncoverenv}<4>
			\draw (1.85,.275) edge[red,thick,out=0,in=135,->,shorten >=.2ex] (3,1);
		\end{uncoverenv}
		
		\draw (.1,.55) -- (.1,.9);
		\draw (.1,.55) -- (1.9,.55);
		\draw (.1,.9) -- (1.9,.9);
		\draw (1.9,.55) -- (1.9,.9);
		\node at (0,.725) [right] {value = ...};
		
		
		\draw (3,-.45) -- (3,1);
		\draw (3,-.45) -- (5,-.45);
		\draw (3,1) -- (5,1);
		\draw (5,-.45) -- (5,1);
		
		\draw (3.1,-.35) -- (3.1,0);
		\draw (3.1,-.35) -- (4.9,-.35);
		\draw (3.1,0) -- (4.9,0);
		\draw (4.9,-.35) -- (4.9,0);
		\node at (3.3,-.225) [right] {prev};
		
		\node[circle, thick, draw=black, minimum size=5pt, inner sep = 0pt] at (3.28,-.175) {};
		\begin{uncoverenv}<-3>
			\draw (3.15,-.175) edge[thick,out=180,in=-45,->,shorten >=.2ex] (2,-.45);
		\end{uncoverenv}
		\begin{uncoverenv}<4>
			\draw (3.15,-.175) edge[red,thick,out=180,in=-45,->,shorten >=.2ex] (2,-.45);
		\end{uncoverenv}
		
		\draw (3.1,.1) -- (3.1,.45);
		\draw (3.1,.1) -- (4.9,.1);
		\draw (3.1,.45) -- (4.9,.45);
		\draw (4.9,.1) -- (4.9,.45);
		\node at (3,.275) [right] {next};
		
		\node[circle, thick, draw=black, minimum size=5pt, inner sep = 0pt] at (4.72,.275) {};
		\draw (4.85,.275) edge[thick,out=0,in=135,->,shorten >=.2ex] (6,1);
		
		
		\draw (3.1,.55) -- (3.1,.9);
		\draw (3.1,.55) -- (4.9,.55);
		\draw (3.1,.9) -- (4.9,.9);
		\draw (4.9,.55) -- (4.9,.9);
		\node at (3,.725) [right] {value = ...};
		
		
		\draw (6,-.45) -- (6,1);
		\draw (6,-.45) -- (8,-.45);
		\draw (6,1) -- (8,1);
		\draw (8,-.45) -- (8,1);
		
		\draw (6.1,-.35) -- (6.1,0);
		\draw (6.1,-.35) -- (7.9,-.35);
		\draw (6.1,0) -- (7.9,0);
		\draw (7.9,-.35) -- (7.9,0);
		\node at (6.3,-.225) [right] {prev};
		
		\node[circle, thick, draw=black, minimum size=5pt, inner sep = 0pt] at (6.28,-.175) {};
		\draw (6.15,-.175) edge[thick,out=180,in=-45,->,shorten >=.2ex] (5,-.45);
		
		\draw (6.1,.1) -- (6.1,.45);
		\draw (6.1,.1) -- (7.9,.1);
		\draw (6.1,.45) -- (7.9,.45);
		\draw (7.9,.1) -- (7.9,.45);
		\node at (6,.275) [right] {next};
		
		\node[circle, thick, draw=black, minimum size=5pt, inner sep = 0pt] at (7.72,.275) {};
		\draw (7.85,.275) edge[thick,out=0,in=135,->,shorten >=.2ex] (9,1);
		
		\draw (6.1,.55) -- (6.1,.9);
		\draw (6.1,.55) -- (7.9,.55);
		\draw (6.1,.9) -- (7.9,.9);
		\draw (7.9,.55) -- (7.9,.9);
		\node at (6,.725) [right] {value = ...};
		
		
		\draw (9,-.45) -- (9,1);
		\draw (9,-.45) -- (11,-.45);
		\draw (9,1) -- (11,1);
		\draw (11,-.45) -- (11,1);
		
		\draw (9.1,-.35) -- (9.1,0);
		\draw (9.1,-.35) -- (10.9,-.35);
		\draw (9.1,0) -- (10.9,0);
		\draw (10.9,-.35) -- (10.9,0);
		\node at (9.3,-.225) [right] {prev};
		
		\node[circle, thick, draw=black, minimum size=5pt, inner sep = 0pt] at (9.28,-.175) {};
		\draw (9.15,-.175) edge[thick,out=180,in=-45,->,shorten >=.2ex] (8,-.45);
		
		\draw (9.1,.1) -- (9.1,.45);
		\draw (9.1,.1) -- (10.9,.1);
		\draw (9.1,.45) -- (10.9,.45);
		\draw (10.9,.1) -- (10.9,.45);
		\node at (9,.275) [right] {next};
		
		
		\node[circle, thick, draw=black, minimum size=5pt, inner sep = 0pt] at (10.72,.275) {};
		\draw (10.85,.275) edge[thick,out=0,in=155,->,shorten >=.2ex] (11.5,.45);
		
		\draw (9.1,.55) -- (9.1,.9);
		\draw (9.1,.55) -- (10.9,.55);
		\draw (9.1,.9) -- (10.9,.9);
		\draw (10.9,.55) -- (10.9,.9);
		\node at (9,.725) [right] {value = ...};

		\draw (11.5,.1) -- (11.5,.45);
		\draw (11.5,.1) -- (12.5,.1);
		\draw (11.5,.45) -- (12.5,.45);
		\draw (12.5,.1) -- (12.5,.45);
		\node at (11.4,.265) [right] {NULL};
		
		
		\begin{uncoverenv}<2->
			\draw (2,-1.5) -- (2,-2.95);
			\draw (2,-1.5) -- (4,-1.5);
			\draw (4,-1.5) -- (4,-2.95);
			\draw (4,-2.95) -- (2,-2.95);
			
			\draw (2.1,-1.6) -- (2.1, -1.95);
			\draw (2.1,-1.6) -- (3.9,-1.6);
			\draw (3.9,-1.6) -- (3.9,-1.95);
			\draw (3.9,-1.95) -- (2.1,-1.95);
			\node at (2,-1.775) [right] {value = ...};
			
			\draw (2.1, -2.05) -- (2.1, -2.4);
			\draw (2.1,-2.05) -- (3.9,-2.05);
			\draw (3.9,-2.05) -- (3.9,-2.4);
			\draw (3.9,-2.4) -- (2.1,-2.4);
			\node at (2,-2.225) [right] {next};
			\node[circle, thick, draw=black, minimum size=5pt, inner sep = 0pt] at (3.72,-2.225) {};
			
			\draw (2.1, -2.5) -- (2.1, -2.85);
			\draw (2.1,-2.5) -- (3.9,-2.5);
			\draw (3.9,-2.5) -- (3.9,-2.85);
			\draw (3.9,-2.85) -- (2.1,-2.85);
			\node at (2.3,-2.725) [right] {prev};
			\node[circle, thick, draw=black, minimum size=5pt, inner sep = 0pt] at (2.28,-2.675) {};
			
			
		\end{uncoverenv}
		\begin{uncoverenv}<3->
			\draw [thick,blue,->,shorten >=.2ex] (3.85,-2.225) to[out=0,in=-35] (4,-1.4) to[out=180-35,in=-45] (2.8,-1.1) to[out=180-45,in=135] (3,1);
			\draw [thick,blue,->,shorten >=.2ex] (2.15,-2.675) to[out=180,in=180+35] (1.8,-1) to[out=35,in=-45] (2,-.45);
		\end{uncoverenv}
		\begin{uncoverenv}<5>
			\draw [thick,red,->,shorten >=.2ex] (1.85,.275) to[out=0,in=90] (2.25,-.5) to[out=270,in=135] (2,-1.5);
			\draw [thick,red,->,shorten >=.2ex] (3.28,-.305) to[out=-90,in=90] (4.6,-1.8) to[out=270,in=-45] (4,-2.95);
		\end{uncoverenv}
		
		
	\end{tikzpicture}

	\begin{itemize}[<+->]
		\item With a pointer in each direction, we can iterate forwards and backwards.
		\item When inserting a value we now have to change a little more.
	\end{itemize}
\end{frame}

%-------------------------------------------------------------------------------------------------

\begin{frame}{Cyclical lists}

	\begin{tikzpicture}[scale=.73,font=\scriptsize]
			
		\draw (-2.5,.1) -- (-1,.1);
		\draw (-2.5,.1) -- (-2.5,.45);
		\draw (-2.5,.45) -- (-1,.45);
		\draw (-1,.45) -- (-1,.1);
		\node[circle, thick, draw=black, minimum size=5pt, inner sep = 0pt] at (-1.18,.275) {};
		\node at (-2.6,.275) [right] {head};
		\draw (-1.05,.275) edge[thick,out=0,in=135,->,shorten >=.2ex] (0,1);
	
		\draw (0,0) -- (0,1);
		\draw (0,0) -- (2,0);
		\draw (0,1) -- (2,1);
		\draw (2,0) -- (2,1);
		
		\draw (.1,.1) -- (.1,.45);
		\draw (.1,.1) -- (1.9,.1);
		\draw (.1,.45) -- (1.9,.45);
		\draw (1.9,.1) -- (1.9,.45);
		\node at (0,.275) [right] {next};
		
		\node[circle, thick, draw=black, minimum size=5pt, inner sep = 0pt] at (1.72,.275) {};
			\draw (1.85,.275) edge[thick,out=0,in=135,->,shorten >=.2ex] (3,1);
		
		\draw (.1,.55) -- (.1,.9);
		\draw (.1,.55) -- (1.9,.55);
		\draw (.1,.9) -- (1.9,.9);
		\draw (1.9,.55) -- (1.9,.9);
		\node at (0,.725) [right] {value = ...};
		
		
		\draw (3,0) -- (3,1);
		\draw (3,0) -- (5,0);
		\draw (3,1) -- (5,1);
		\draw (5,0) -- (5,1);
		
		\draw (3.1,.1) -- (3.1,.45);
		\draw (3.1,.1) -- (4.9,.1);
		\draw (3.1,.45) -- (4.9,.45);
		\draw (4.9,.1) -- (4.9,.45);
		\node at (3,.275) [right] {next};
		
		\node[circle, thick, draw=black, minimum size=5pt, inner sep = 0pt] at (4.72,.275) {};
		\draw (4.85,.275) edge[thick,out=0,in=135,->,shorten >=.2ex] (6,1);
		
		
		\draw (3.1,.55) -- (3.1,.9);
		\draw (3.1,.55) -- (4.9,.55);
		\draw (3.1,.9) -- (4.9,.9);
		\draw (4.9,.55) -- (4.9,.9);
		\node at (3,.725) [right] {value = ...};
		
		
		\draw (6,0) -- (6,1);
		\draw (6,0) -- (8,0);
		\draw (6,1) -- (8,1);
		\draw (8,0) -- (8,1);
		
		\draw (6.1,.1) -- (6.1,.45);
		\draw (6.1,.1) -- (7.9,.1);
		\draw (6.1,.45) -- (7.9,.45);
		\draw (7.9,.1) -- (7.9,.45);
		\node at (6,.275) [right] {next};
		
		\node[circle, thick, draw=black, minimum size=5pt, inner sep = 0pt] at (7.72,.275) {};
		\draw (7.85,.275) edge[thick,out=0,in=135,->,shorten >=.2ex] (9,1);
		
		\draw (6.1,.55) -- (6.1,.9);
		\draw (6.1,.55) -- (7.9,.55);
		\draw (6.1,.9) -- (7.9,.9);
		\draw (7.9,.55) -- (7.9,.9);
		\node at (6,.725) [right] {value = ...};
		
		
		\draw (9,0) -- (9,1);
		\draw (9,0) -- (11,0);
		\draw (9,1) -- (11,1);
		\draw (11,0) -- (11,1);
		
		\draw (9.1,.1) -- (9.1,.45);
		\draw (9.1,.1) -- (10.9,.1);
		\draw (9.1,.45) -- (10.9,.45);
		\draw (10.9,.1) -- (10.9,.45);
		\node at (9,.275) [right] {next};
		
		
		\node[circle, thick, draw=black, minimum size=5pt, inner sep = 0pt] at (10.72,.275) {};
		\draw [thick,->,shorten >=.2ex] (10.72,.175) to[out=-90,in=0] (9,-.5) to[out=180,in=0] (3.5,-.5) to[out=180,in=-45] (2,0);
		
		\draw (9.1,.55) -- (9.1,.9);
		\draw (9.1,.55) -- (10.9,.55);
		\draw (9.1,.9) -- (10.9,.9);
		\draw (10.9,.55) -- (10.9,.9);
		\node at (9,.725) [right] {value = ...};
		
	\end{tikzpicture}\\\ \\\ \\\ \\
	\pause
	\begin{itemize}
		\item You can also implement \textit{double-linked cyclical lists}.
	\end{itemize}
\end{frame}

%-------------------------------------------------------------------------------------------------

\begin{frame}{Binary trees}
	\begin{columns}[T]
		\column{.5\textwidth}
		\begin{tikzpicture}[scale=.12,font=\scriptsize]
			\draw (-5,0) -- (5,0);
			\draw (5,0) -- (5,-2);
			\draw (5,-2) -- (-5,-2);
			\draw (-5,-2) -- (-5,0);
			
			\node at (-1,.7)[below] {root};
			\node[circle, thick, draw=black, minimum size=4pt, inner sep = 0pt] at (4,-1) {};
			\draw [thick,->,shorten >=.2ex,shorten <=.4ex] (4,-1) to[out=-90,in=90] (0,-7);
			
			
			\draw (-5,-7) -- (5,-7);
			\draw (5,-7) -- (5,-15);
			\draw (5,-15) -- (-5,-15);
			\draw (-5,-15) -- (-5,-7);
			
			\draw (-4.5,-7.5) -- (4.5,-7.5);
			\draw (4.5,-7.5) -- (4.5,-9.5);
			\draw (4.5,-9.5) -- (-4.5,-9.5);
			\draw (-4.5,-9.5) -- (-4.5,-7.5);
			\node at (0,-6.7)[below] {value};
			
			\draw (-4.5,-10) -- (4.5,-10);
			\draw (4.5,-10) -- (4.5,-12);
			\draw (4.5,-12) -- (-4.5,-12);
			\draw (-4.5,-12) -- (-4.5,-10);
			\node at (-1,-9.1)[below] {right};
			\node[circle, thick, draw=black, minimum size=4pt, inner sep = 0pt] at (3.5,-11) {};
			\draw [thick,->,shorten >=.2ex,shorten <=.4ex] (3.5,-11) to[out=0,in=90] (10,-22);
			
			\draw (-4.5,-12.5) -- (4.5,-12.5);
			\draw (4.5,-12.5) -- (4.5,-14.5);
			\draw (4.5,-14.5) -- (-4.5,-14.5);
			\draw (-4.5,-14.5) -- (-4.5,-12.5);
			\node at (1,-11.7)[below] {left};
			\node[circle, thick, draw=black, minimum size=4pt, inner sep = 0pt] at (-3.5,-13.5) {};
			\draw [thick,->,shorten >=.2ex,shorten <=.4ex] (-3.5,-13.5) to[out=-180,in=90] (-10,-22);
			
			
			\draw (-15,-22) -- (-5,-22);
			\draw (-5,-22) -- (-5,-30);
			\draw (-5,-30) -- (-15,-30);
			\draw (-15,-30) -- (-15,-22);
			
			\draw (-14.5,-22.5) -- (-5.5,-22.5);
			\draw (-5.5,-22.5) -- (-5.5,-24.5);
			\draw (-5.5,-24.5) -- (-14.5,-24.5);
			\draw (-14.5,-24.5) -- (-14.5,-22.5);
			\node at (-10,-21.7)[below] {value};
			
			\draw (-14.5,-25) -- (-5.5,-25);
			\draw (-5.5,-25) -- (-5.5,-27);
			\draw (-5.5,-27) -- (-14.5,-27);
			\draw (-14.5,-27) -- (-14.5,-25);
			\node at (-11,-24.1)[below] {right};
			\node[circle, thick, draw=black, minimum size=4pt, inner sep = 0pt] at (-6.5,-26) {};
			\draw [thick,->,shorten >=.2ex,shorten <=.4ex] (-6.5,-26) to[out=0,in=90] (-3,-37);
			
			\draw (-14.5,-27.5) -- (-5.5,-27.5);
			\draw (-5.5,-27.5) -- (-5.5,-29.5);
			\draw (-5.5,-29.5) -- (-14.5,-29.5);
			\draw (-14.5,-29.5) -- (-14.5,-27.5);
			\node at (-9,-26.7)[below] {left};
			\node[circle, thick, draw=black, minimum size=4pt, inner sep = 0pt] at (-13.5,-28.5) {};
			\draw [thick,->,shorten >=.2ex,shorten <=.4ex] (-13.5,-28.5) to[out=180,in=90] (-20,-37);
			
			
			
			\draw (5,-22) -- (15,-22);
			\draw (15,-22) -- (15,-30);
			\draw (15,-30) -- (5,-30);
			\draw (5,-30) -- (5,-22);
			
			\draw (5.5,-22.5) -- (14.5,-22.5);
			\draw (14.5,-22.5) -- (14.5,-24.5);
			\draw (14.5,-24.5) -- (5.5,-24.5);
			\draw (5.5,-24.5) -- (5.5,-22.5);
			\node at (10,-21.7)[below] {value};
			
			\draw (5.5,-25) -- (14.5,-25);
			\draw (14.5,-25) -- (14.5,-27);
			\draw (14.5,-27) -- (5.5,-27);
			\draw (5.5,-27) -- (5.5,-25);
			\node at (9,-24.1)[below] {right};
			\node[circle, thick, draw=black, minimum size=4pt, inner sep = 0pt] at (13.5,-26) {};
			\draw [thick,->,shorten >=.2ex,shorten <=.4ex] (13.5,-26) to[out=0,in=90] (20,-37);
			
			\draw (5.5,-27.5) -- (14.5,-27.5);
			\draw (14.5,-27.5) -- (14.5,-29.5);
			\draw (14.5,-29.5) -- (5.5,-29.5);
			\draw (5.5,-29.5) -- (5.5,-27.5);
			\node at (11,-26.7)[below] {left};
			\node[circle, thick, draw=black, minimum size=4pt, inner sep = 0pt] at (6.5,-28.5) {};
			\draw [thick,->,shorten >=.2ex,shorten <=.4ex] (6.5,-28.5) to[out=180,in=90] (3,-37);
			
			
			
			
			\draw (-25,-37) -- (-15,-37);
			\draw (-15,-37) -- (-15,-45);
			\draw (-15,-45) -- (-25,-45);
			\draw (-25,-45) -- (-25,-37);
			
			\draw (-24.5,-37.5) -- (-15.5,-37.5);
			\draw (-15.5,-37.5) -- (-15.5,-39.5);
			\draw (-15.5,-39.5) -- (-24.5,-39.5);
			\draw (-24.5,-39.5) -- (-24.5,-37.5);
			\node at (-20,-36.7)[below] {value};
			
			\draw (-24.5,-40) -- (-15.5,-40);
			\draw (-15.5,-40) -- (-15.5,-42);
			\draw (-15.5,-42) -- (-24.5,-42);
			\draw (-24.5,-42) -- (-24.5,-40);
			\node at (-21,-39.1)[below] {right};
			\node[circle, thick, draw=black, minimum size=4pt, inner sep = 0pt] at (-16.5,-41) {};
			\draw [thick,->,shorten >=.2ex,shorten <=.4ex] (-16.5,-41) to[out=0,in=90] (-10,-52);
			
			\draw (-24.5,-42.5) -- (-15.5,-42.5);
			\draw (-15.5,-42.5) -- (-15.5,-44.5);
			\draw (-15.5,-44.5) -- (-24.5,-44.5);
			\draw (-24.5,-44.5) -- (-24.5,-42.5);
			\node at (-19,-41.7)[below] {left};
			\node[circle, thick, draw=black, minimum size=4pt, inner sep = 0pt] at (-23.5,-43.5) {};
			\draw [thick,->,shorten >=.2ex,shorten <=.4ex] (-23.5,-43.5) to[out=-180,in=90] (-30,-52);
	
	
	
	
			\draw (0,-37) -- (10,-37);
			\draw (10,-37) -- (10,-45);
			\draw (10,-45) -- (0,-45);
			\draw (0,-45) -- (0,-37);
			
			\draw (.5,-37.5) -- (9.5,-37.5);
			\draw (9.5,-37.5) -- (9.5,-39.5);
			\draw (9.5,-39.5) -- (.5,-39.5);
			\draw (.5,-39.5) -- (.5,-37.5);
			\node at (5,-36.7)[below] {value};
			
			\draw (.5,-40) -- (9.5,-40);
			\draw (9.5,-40) -- (9.5,-42);
			\draw (9.5,-42) -- (.5,-42);
			\draw (.5,-42) -- (.5,-40);
			\node at (4,-39.1)[below] {right};
			\node[circle, thick, draw=black, minimum size=4pt, inner sep = 0pt] at (8.5,-41) {};
			\draw [thick,->,shorten >=.2ex,shorten <=.4ex] (8.5,-41) to[out=0,in=90] (15,-52);
			
			\draw (.5,-42.5) -- (9.5,-42.5);
			\draw (9.5,-42.5) -- (9.5,-44.5);
			\draw (9.5,-44.5) -- (.5,-44.5);
			\draw (.5,-44.5) -- (.5,-42.5);
			\node at (6,-41.7)[below] {left};
			\node[circle, thick, draw=black, minimum size=4pt, inner sep = 0pt] at (1.5,-43.5) {};
			\draw [thick,->,shorten >=.2ex,shorten <=.4ex] (1.5,-43.5) to[out=-180,in=90] (-3,-52);
	
	
	
			\draw (-35,-52) -- (-25,-52);
			\draw (-35,-52) -- (-35,-54);
			\draw (-25,-52) -- (-25,-54);
			\draw (-35,-54) -- (-25,-54);
			\node at (-30,-51.2)[below] {NULL};
	
	
	
	
			\draw (-18,-52) -- (-8,-52);
			\draw (-18,-52) -- (-18,-54);
			\draw (-8,-52) -- (-8,-54);
			\draw (-18,-54) -- (-8,-54);
			\node at (-13,-51.2)[below] {NULL};
	
	
	
	
			\draw (-5,-52) -- (5,-52);
			\draw (-5,-52) -- (-5,-54);
			\draw (5,-52) -- (5,-54);
			\draw (-5,-54) -- (5,-54);
			\node at (0,-51.2)[below] {NULL};
	
	
	
	
			\draw (10,-52) -- (20,-52);
			\draw (20,-52) -- (20,-54);
			\draw (10,-52) -- (10,-54);
			\draw (10,-54) -- (20,-54);
			\node at (15,-51.2)[below] {NULL};
	
	
	
			\draw (-12,-37) -- (-2,-37);
			\draw (-12,-37) -- (-12,-39);
			\draw (-2,-37) -- (-2,-39);
			\draw (-12,-39) -- (-2,-39);
			\node at (-7,-36.2)[below] {NULL};
	
	
	
	
			\draw (15,-37) -- (25,-37);
			\draw (15,-37) -- (15,-39);
			\draw (25,-37) -- (25,-39);
			\draw (15,-39) -- (25,-39);
			\node at (20,-36.2)[below] {NULL};
		\end{tikzpicture}
		\column{.6\textwidth}
		\begin{itemize}
			\item To minimize the complexity of list algorithms, binary trees are constructed
			\item Such algorithms are e.g. searching, sorting, adding
		\end{itemize}
	\end{columns}
\end{frame}

%-------------------------------------------------------------------------------------------------
\section{The Dungeon}
%-------------------------------------------------------------------------------------------------

\begin{frame}{Back to the Dungeon}
	Again, we prepared a version of the ASCII-Dungeon for you in the repository (\url{https://github.com/fsr/c-slides}) in folder \textit{materials/2\_before/}.\\
	\ \\
	
	\begin{itemize}
		\item There are two monsters in the game, separated from each other.
		\item Edit the code to have a double chained list of monsters.
		\item You may want to add a few functions:
		\begin{itemize}
			\item init\_monster\_list
			\item free\_monster\_list
			\item add\_monster
			\item print\_monster\_list
		\end{itemize}
	\end{itemize}

\end{frame}

%-------------------------------------------------------------------------------------------------
\section{Pointers to functions}
%-------------------------------------------------------------------------------------------------

\begin{frame}[fragile]{Higher order functions}
	It is possible to pass a function as an argument of another function.\\
	Doing so, you get a function that can call different other functions.\\
	\bigskip
	Since functions are placed in the text segment of you program, you simply pass the address of the function as a so called function pointer:\\
	\begin{lstlisting}
<return type> (*<function name>)(<parameter list>);
\end{lstlisting}
	\bigskip
	Example:
	\begin{lstlisting}
int (*op)(int a, int b) = add; /* Defined somewhere else */
int (*print)(const char *) = puts;	/* From stdlib */
\end{lstlisting}
	Note that the $\&$ operator can be omitted.
\end{frame}

%-------------------------------------------------------------------------------------------------

\begin{frame}[fragile]{Example}
	\begin{lstlisting}
#include <stdio.h>

int add(int a, int b) {
	return a + b;
}
int sub(int a, int b) {
	return a - b;
}
void printFunc(int (*f)(int, int), int a, int b) {
	 printf("%d\n",f(a, b));
}
int main(void) {
	printFunc(add, 1, 2);
	printFunc(sub, 1, 2);
	return 0;
}
\end{lstlisting}

\end{frame}

%-------------------------------------------------------------------------------------------------

\begin{frame}{Mapping}
	Pointers to functions are often used for mapping. If you want to iterate through a list and call a function for every list item, you could do this with a new loop each time, but you also could write a function that takes a list and a function.\\
	\ \\
	If only we had an example to try this out...\\
	Oh wait... we do: The Dungeon.\\\ \\
	\begin{itemize}
		\item Write a mapping function that takes a list of monsters and a function.
		\item Edit the print\_monster\_list function to a print\_entity function.
		\item Now you can print the monsters list by calling the mapping function passing the list and the print\_entity function.
	\end{itemize}
\end{frame}

%-------------------------------------------------------------------------------------------------
\end{document}
